\section{Threat model and trusted base}
\label{sec:threat}

The kernel sits between an agent and the tool servers whose transport it owns:
child processes, fronted HTTP tools, proxied sidecars. It does not confine the
agent, whose model, orchestrator, prompt, and libraries are untrusted; one that
opens its own socket reaches whatever that socket reaches.

The adversary is given everything outside the receiving kernel: every field
this paper names, the peer kernel and its signing key, the network, where
it may drop, reorder, and replay, and any receipt, statement, lease, or
governance record it has observed. It holds neither the receiving kernel's
signing key nor write access to its stores. The question is what it can cause a
receiver to dispatch.

The trusted base is the receiving kernel's process, its signing key, and its
stores; everything outside is verified or assumed; the assumptions appear once,
in Table~\ref{tab:assumptions}.

\scriptsize
\begin{longtable}{@{}>{\raggedright\arraybackslash}p{4.3cm}p{10.9cm}@{}}
\caption{Assumptions outside the receiving kernel, by the identifiers of the implementation's assumption registry. The registry has \PSAssumptionCount{} entries; the ten this protocol rests on are listed.}\label{tab:assumptions}\\
\toprule
Identifier & What is assumed \\
\midrule
\endfirsthead
\toprule
Identifier & What is assumed \\
\midrule
\endhead
\bottomrule
\endfoot
ASSUME-ED25519 & Ed25519 meets existential unforgeability under chosen-message attack for keys the kernel trusts. \\
ASSUME-SHA256 & SHA-256 is collision resistant over receipt and statement preimages. \\
ASSUME-CANONICAL-JSON & The production canonical encoder agrees byte for byte with the mechanized one on strings, bounded integers, finite arrays, and objects ordered by UTF-16 code unit. Floats are outside that domain. \\
ASSUME-SQLITE-ATOMICITY & A single-row insert commits atomically, so a primary-key conflict is a reliable witness that a row existed. Cross-row recovery and ordering are not assumed. \\
ASSUME-OS-CLOCK & The injected clock reports Unix time within the operator's declared tolerance. \\
ASSUME-NETWORK-TRANSPORT, ASSUME-TLS & Delivery is not assumed reliable; arriving messages are assumed not rewritten below the signature checks, and remote surfaces authenticate their endpoints. \\
ASSUME-SUBPROCESS-ISOLATION & Effects after an allow, and isolation between tool processes, rest on operating-system boundaries. \\
ASSUME-FEDERATED-ORIGIN-CLASSIFICATION & The receiving edge marks cross-organization origin from a transport-authenticated peer identity; a request misclassified as local never reaches these checks. \\
ASSUME-GOSSIP-FAIRNESS-PARTITION-BOUND & Revocation liveness only: recurring delivery under weak fairness, bounded skew, operator-bounded partition healing. Safety does not depend on it. \\
\end{longtable}
\normalsize

Two assumptions are deliberately absent. The first is that the two signing keys
on a co-signed statement belong to different organizations: one operator holding
both satisfies every check here. The second is domain separation: one passport
key signs three preimage families, the receipt signing body, the co-signing
body, and a statement's pre-authentication encoding, so ASSUME-ED25519 is an
assumption about a key with that reuse. A pinned peer obtains this kernel's
signature over any bytes that rebuild as one of the last two, so a
countersignature fixes a record's shape and parties, not its truth.
